%! Vector-Valued Functions — Unit 4, Lesson 4.9 — Homework (Answer Key)
\documentclass[10pt]{article}
\usepackage{precalculus-article}
\usepackage{precalculus-key}   % pulls in -boxes; do NOT also load -boxes
\newcommand{\TallMath}[1]{$\displaystyle #1\rule[-1.4em]{0pt}{3.2em}$}

\begin{document}

\pageheader{Unit 4, Lesson 4.9}{Homework --- Answer Key}
\namedateperiod

\vspace{0.12in}

\begin{notesbox}{Vector-Valued Functions Practice --- Key}
\small

\textbf{1.}\quad $f(t)=(5t-2,\;t^{2}+1)$.

\begin{enumerate}[label=(\alph*)]
  \item $\vec{p}(t) = $ \ans{$\langle 5t-2,\;t^{2}+1\rangle$}
  \item $\vec{p}(3) = \langle 5(3)-2,\;9+1\rangle = $ \ans{$\langle 13,\;10\rangle$}
\end{enumerate}

\vspace{0.12in}

\textbf{2.}\quad $\|\vec{p}(3)\| = \sqrt{13^{2}+10^{2}} = \sqrt{169+100}
= \sqrt{269} \approx $ \ans{$16.4$}

\vspace{0.12in}

\textbf{3.}\quad $\vec{v}=\langle -3,\;5\rangle$.

\begin{enumerate}[label=(\alph*)]
  \item \ans{Moving left, because $x'=-3<0$.}
  \item \ans{Moving up, because $y'=5>0$.}
\end{enumerate}

\vspace{0.12in}

\textbf{4.}\quad $\vec{v}(2) = \langle 6,\;-8\rangle$ \qquad
Speed $= \sqrt{36+64} = \sqrt{100} = $ \ans{$10$}

\vspace{0.12in}

\textbf{5.}\quad $\vec{p}(t)=\langle 4\cos t,\;3\sin t\rangle$.

$\|\vec{p}(t)\|^{2} = 16\cos^{2}t + 9\sin^{2}t
= 9(\cos^{2}t+\sin^{2}t) + 7\cos^{2}t
= 9 + 7\cos^{2}t.$

This is maximized when $\cos^{2}t=1$, i.e.\ when $t=0$ or $t=\pi$.
At $t=0$: $\|\vec{p}(0)\|=\sqrt{9+7}=4$. \quad
At $t=\pi/2$: $\|\vec{p}(\pi/2)\|=\sqrt{9+0}=3$.

The maximum distance is $4$, achieved at $t=0$ (and $t=\pi$).

\ans{(A) $t=0$}

\begin{teachernote}
  Many students guess $t=\pi/2$ because $\sin$ reaches its maximum there.
  Walk through the algebraic argument: $\|\vec{p}(t)\|^{2}=9+7\cos^{2}t$
  is maximized by maximizing $\cos^{2}t$, not $\sin^{2}t$.  This problem
  builds reasoning skill 3.C (supporting conclusions with data/rationale).
\end{teachernote}

\end{notesbox}

\end{document}
